\documentclass[
    kokkoshighlight,
    % print,
]{./cheat_sheet_kokkos}

\usepackage[T1]{fontenc}
\usepackage[utf8]{inputenc}

% document properties
\title{Utilization cheat sheet for Kokkos}
\author{CExA}
\date{\today}

\begin{document}
\maketitle

\begin{multicols}{2}

\section{Header}

\begin{minted}{c++}
#include <Kokkos_Core.hpp>
\end{minted}

\section{Initialization}

\subsection{Initialize and finalize}

\begin{minted}{c++}
int main(int argc, char* argv[]) {
    Kokkos::initialize(argc, argv);
    { /* ... */ }
    Kokkos::finalize();
}
\end{minted}

\subsection{Scope guard}

\begin{minted}{c++}
int main(int argc, char* argv[]) {
    Kokkos::ScopeGuard kokkos(argc, argv);
    /* ... */
}
\end{minted}

\section{Kokkos concepts}

\subsection{Execution spaces}

\begin{tblr}[theme=kokkostable]{Xll}
    Execution space & Device backend & Host backend \\
    \mintinline{c++}{Kokkos::DefaultExecutionSpace} & On device & On host \\
    \mintinline{c++}{Kokkos::DefaultHostExecutionSpace} & On host & On host \\
\end{tblr}

\subsection{Memory spaces}

\subsubsection{Generic memory spaces}

\begin{tblr}[theme=kokkostable]{lXX}
    Memory space & Device backend & Host backend \\
    \mintinline{c++}{Kokkos::DefaultExecutionSpace::memory_space} & On dev. & On host \\
    \mintinline{c++}{Kokkos::DefaultHostExecutionSpace::memory_space} & On host & On host \\
\end{tblr}

\subsubsection{Specific memory spaces}

\begin{tblr}[theme=kokkostable]{lX}
    Memory space & Description \\
    \mintinline{c++}{Kokkos::HostSpace} & Accessible from the host but maybe not from the device \\
    \mintinline{c++}{Kokkos::SharedSpace} & Accessible from the host and the device; copy managed by the driver \\
    \mintinline{c++}{Kokkos::SharedHostPinnedSpace} & Accessible from the host and the device; zero copy access in small chunks \\
\end{tblr}

\section{Memory management}

\subsection{View}

\subsubsection{Create}

\begin{minted}{c++}
Kokkos::View<DataType, LayoutType, MemorySpace, MemoryTraits> view("label", numberOfElementsAtRuntimeI, numberOfElementsAtRuntimeJ);
\end{minted}

\begin{tblr}[theme=kokkostable]{lX}
    Template arg. & Description \\
    \mintinline{c++}{DataType} & \mintinline{c++}{ScalarType} for the data type, followed by a \mintinline{c++}{*} for each runtime dimension, then by a \mintinline{c++}{[numberOfElements]} for each compile time dimension, mandatory \\
    \mintinline{c++}{LayoutType} & See memory layouts, optional \\
    \mintinline{c++}{MemorySpace} & See memory spaces, optional \\
    \mintinline{c++}{MemoryTraits} & See memory traits, optional \\
\end{tblr}

The order of template arguments is important.

\subsubsection{Manage}

\begin{tblr}[theme=kokkostable]{lX}
    Method & Description \\
    \mintinline{c++}{(i, j...)} & Returns and sets the value at index \mintinline{c++}{i},
    \mintinline{c++}{j}, etc. \\
    \mintinline{c++}{size()} & Returns the total number of elements in the view \\
    \mintinline{c++}{rank()} & Returns the number of dimensions \\
    \mintinline{c++}{layout()} & Returns the layout of the view \\
    \mintinline{c++}{extent(dim)} & Returns the number of elements in the requested dimension \\
    \mintinline{c++}{data()} & Returns a pointer to the underlying data \\
\end{tblr}

Resize and preserve content

\begin{minted}{c++}
Kokkos::resize(view, newNumberOfElementsI, newNumberOfElementsJ...);
\end{minted}

Reallocate and do not preserve content

\begin{minted}{c++}
Kokkos::realloc(view, newNumberOfElementsI, newNumberOfElementsJ...);
\end{minted}

\subsection{Memory Layouts}

\begin{tblr}[theme=kokkostable]{lXl}
    Layout & Description & Default \\
    \mintinline{c++}{Kokkos::LayoutRight} & Strides increase from the right most to the left most dimension, also known as row-major or C-like & CPU \\
    \mintinline{c++}{Kokkos::LayoutLeft} & Strides increase from the left most to the right most dimension, also known as column-major or Fortran-like & GPU \\
    \mintinline{c++}{Kokkos::LayoutStride} & Strides can be arbitrary for each dimension &  \\
\end{tblr}

By default, a layout suited for loops on the high frequency index is used.

\subsection{Memory trait}

Memory traits are indicated with \mintinline{c++}{Kokkos::MemoryTraits<>} and are combined with the \mintinline{c++}{|} (pipe) operator.

\begin{tblr}[theme=kokkostable]{lX}
    Memory trait & Description \\
    \mintinline{c++}{Kokkos::Unmanaged} & The allocation has to be managed manually \\
    \mintinline{c++}{Kokkos::Atomic} & All accesses to the view are atomic \\
    \mintinline{c++}{Kokkos::RandomAccess} & Hint that the view is used in a random access manner; if the view is also \mintinline{c++}{const} this may trigger more efficient load operations on GPUs \\
    \mintinline{c++}{Kokkos::Restrict} & There is no aliasing of the view by other data structures in the current scope \\
\end{tblr}

\subsection{Deep copy}

\begin{minted}{c++}
Kokkos::deep_copy(destination, source);
\end{minted}

The views must have the same dimensions, data type, and reside in the same memory space (mirror views can be deep copied on different memory spaces).

\subsection{Mirror view}

Create and always allocate on host

\begin{minted}{c++}
auto mirrorView = Kokkos::create_mirror(view);
\end{minted}

Create and allocate on host if source view is not in host space

\begin{minted}{c++}
auto mirrorView = Kokkos::create_mirror_view(view);
\end{minted}

Create, allocate and synchronize if source view is not in same space as destination view

\begin{minted}{c++}
auto mirrorView = Kokkos::create_mirror_view_and_copy(ExecutionSpace(), view);
\end{minted}

\subsection{Subview}

A subview has the same reference count as its parent view, so the parent view won't be deallocated before all subviews go away.

\begin{minted}{c++}
auto subview = Kokkos::subview(view, selector1, selector2, ...);
\end{minted}

\begin{tblr}[theme=kokkostable]{lX}
    Subset selector & Description \\
    \mintinline{c++}{Kokkos::ALL} & All elements in this dimension \\
    \mintinline{c++}{Kokkos::pair(first, last)} & Range of elements in this dimension \\
    \mintinline{c++}{value} & Specific element in this dimension \\
\end{tblr}

\subsection{Scatter view (experimental)}

\subsubsection{Specific header}

\begin{minted}{c++}
#include <Kokkos_ScatterView.hpp>
\end{minted}

\subsubsection{Create}

\begin{minted}{c++}
auto scatterView = Kokkos::Experimental::create_scatter_view<Operation, Duplication, Contribution>(targetView);
\end{minted}

\begin{tblr}[theme=kokkostable]{lX}
    Template arg. & Description \\
    \mintinline{c++}{Operation} & See scatter operation; defaults to \mintinline{c++}{Kokkos::Experimental::ScatterSum} \\
    \mintinline{c++}{Duplication} & Whether to duplicate the grid or not; choices are
    \mintinline{c++}{Kokkos::Experimental::ScatterDuplicated}, and
    \mintinline{c++}{Kokkos::Experimental::ScatterNonDuplicated}; defaults to the option that is the most optimised for \mintinline{c++}{targetView}'s execution space \\
    \mintinline{c++}{Contribution} & Whether to contribute using atomics or not; choices are \mintinline{c++}{Kokkos::Experimental::ScatterAtomic}, or
    \mintinline{c++}{Kokkos::Experimental::ScatterNonAtomic}; defaults to the option that is the most optimised for \mintinline{c++}{targetView}'s execution space \\
\end{tblr}

\subsubsection{Scatter operation}

\begin{tblr}[theme=kokkostable]{Xl}
    Operation & Description \\
    \mintinline{c++}{Kokkos::Experimental::ScatterSum} & Sum \\
    \mintinline{c++}{Kokkos::Experimental::ScatterProd} & Product \\
    \mintinline{c++}{Kokkos::Experimental::ScatterMin} & Minimum value \\
    \mintinline{c++}{Kokkos::Experimental::ScatterMax} & Maximum value \\
\end{tblr}

\subsubsection{Scatter, compute, and gather}

\begin{minted}{c++}
Kokkos::parallel_for(
    "label",
    /* ... */,
    KOKKOS_LAMBDA (/* ... */) {
        // scatter
        auto scatterAccess = scatterView.access();

        // compute
        scatterAccess(/* index */) /* operation */ /* contribution */;
    }
);

// gather
Kokkos::Experimental::contribute(targetView, scatterView);
\end{minted}

\section{Parallel constructs}

\subsection{For loop}

\begin{minted}{c++}
Kokkos::parallel_for(
    "label",
    ExecutionPolicy</* ... */>(/* ... */),
    KOKKOS_LAMBDA (/* ... */) { /* ... */ }
);
\end{minted}

\subsection{Reduction}

\begin{minted}{c++}
ScalarType result;
Kokkos::parallel_reduce(
    "label",
    ExecutionPolicy</* ... */>(/* ... */),
    KOKKOS_LAMBDA (/* ... */, ScalarType& resultLocal) { /* ... */ },
    Kokkos::ReducerConcept<ScalarType>(result)
);
\end{minted}

With \mintinline{c++}{Kokkos::ReducerConcept} being one of the following:

\begin{tblr}[theme=kokkostable]{llX}
    Reducer & Operation & Description \\
    \mintinline{c++}{Kokkos::BAnd} & \mintinline{c++}{&} & Binary and \\
    \mintinline{c++}{Kokkos::BOr} & \mintinline{c++}{|} & Binary or \\
    \mintinline{c++}{Kokkos::LAnd} & \mintinline{c++}{&&} & Logical and \\
    \mintinline{c++}{Kokkos::LOr} & \mintinline{c++}{||} & Logical or \\
    \mintinline{c++}{Kokkos::Max} & \mintinline{c++}{std::max} & Maximum \\
    \mintinline{c++}{Kokkos::MaxLoc} & \mintinline{c++}{std::max_element} & Maximum and associated index \\
    \mintinline{c++}{Kokkos::Min} & \mintinline{c++}{std::min} & Minimum \\
    \mintinline{c++}{Kokkos::MinLoc} & \mintinline{c++}{std::min_element} & Minimum and associated index \\
    \mintinline{c++}{Kokkos::MinMax} & \mintinline{c++}{std::minmax} & Minimum and maximum \\
    \mintinline{c++}{Kokkos::MinMaxLoc} & \mintinline{c++}{std::minmax_element} & Minimum and maximum and associated indices \\
    \mintinline{c++}{Kokkos::Prod} & \mintinline{c++}{*} & Product \\
    \mintinline{c++}{Kokkos::Sum} & \mintinline{c++}{+} & Sum \\
\end{tblr}

A scalar value may be passed, for which the reduction is limited to a sum.
When using the \mintinline{c++}{TeamVectorMDRange}, the \mintinline{c++}{TeamThreadMDRange}, or the \mintinline{c++}{ThreadVectorMDRange} execution policy, only a scalar value may be passed, for which the reduction is also limited to a sum.

\subsection{Fences}

\subsubsection{Global fence}

\begin{minted}{c++}
Kokkos::fence("label");
\end{minted}

\subsubsection{Execution space fence}

\begin{minted}{c++}
ExecutionSpace().fence("label");
\end{minted}

\subsubsection{Team barrier}

\begin{minted}{c++}
Kokkos::TeamPolicy<>::member_type().team_barrier();
\end{minted}

\section{Execution policy}

\subsection{Create}

\begin{minted}{c++}
ExecutionPolicy<ExecutionSpace, Schedule, IndexType, LaunchBounds, WorkTag> policy(/* ... */);
\end{minted}

\begin{tblr}[theme=kokkostable]{lX}
    Template arg. & Description \\
    \mintinline{c++}{ExecutionSpace} & See execution spaces; defaults to \mintinline{c++}{Kokkos::DefaultExecutionSpace} \\
    \mintinline{c++}{Schedule} & How to schedule work items; defaults to machine and backend specifics \\
    \mintinline{c++}{IndexType} & Integer type to be used for the index; defaults to
    \mintinline{c++}{int64_t} \\
    \mintinline{c++}{LaunchBounds} & Hints for CUDA and HIP launch bounds \\
    \mintinline{c++}{WorkTag} & Empty tag class to call the functor \\
\end{tblr}

\subsection{Ranges}

\subsubsection{One-dimensional range}

\begin{minted}{c++}
Kokkos::RangePolicy<ExecutionSpace, Schedule, IndexType LaunchBounds, WorkTag> policy(first, last);
\end{minted}

If the range starts at 0 and uses default parameters, can be replaced by just the number of elements.

\subsubsection{Multi-dimensional (dimension 2)}

\begin{minted}{c++}
Kokkos::MDRangePolicy<ExecutionSpace, Schedule, IndexType, LaunchBounds, WorkTag, Kokkos::Rank<2>> policy({firstI, firstJ}, {lastI, lastJ});
\end{minted}

\subsection{Hierarchical parallelism}

\subsubsection{Team policy}

\begin{minted}{c++}
Kokkos::TeamPolicy<ExecutionSpace, Schedule, IndexType, LaunchBounds, WorkTag> policy(leagueSize, teamSize);
\end{minted}

Usually, \mintinline{c++}{teamSize} is replaced by \mintinline{c++}{Kokkos::AUTO} to let Kokkos determine it. A kernel running in a team policy has a \mintinline{c++}{Kokkos::TeamPolicy<>::member_type} argument:

\begin{tblr}[theme=kokkostable]{lX}
    Method & Description \\
    \mintinline{c++}{league_size()} & Number of teams in the league \\
    \mintinline{c++}{league_rank()} & Index of the team within the league \\
    \mintinline{c++}{team_size()} & Number of threads in the team \\
    \mintinline{c++}{team_rank()} & Index of the thread within the team \\
\end{tblr}

Note that nested parallel constructs do not use \mintinline{c++}{KOKKOS_LAMBDA} to create lambdas. One must use the C++ syntax, for example \mintinline{c++}{[=]} or \mintinline{c++}{[&]}.

\subsubsection{Team vector level (2-level hierarchy)}

\begin{minted}{c++}
Kokkos::parallel_for(
    "label",
    Kokkos::TeamPolicy(numberOfElementsI, Kokkos::AUTO),
    KOKKOS_LAMBDA (const Kokkos::TeamPolicy<>::member_type& teamMember) {
        const int i = teamMember.team_rank();

        Kokkos::parallel_for(
            Kokkos::TeamVectorRange(teamMember, firstJ, lastJ),
            [=] (const int j) { /* ... */ }
        );
    }
);
\end{minted}

\paragraph{One-dimensional range}

\begin{minted}{c++}
Kokkos::TeamVectorRange range(teamMember, firstJ, lastJ);
\end{minted}

\paragraph{Multi-dimensional range (dimension 2)}

\begin{minted}{c++}
Kokkos::TeamVectorMDRange<Kokkos::Rank<2>, Kokkos::TeamPolicy<>::member_type> range(teamMember, numberOfElementsJ, numberOfElementsK);
\end{minted}

\subsubsection{Team thread vector level (3-level hierarchy)}

\begin{minted}{c++}
Kokkos::parallel_for(
    "label",
    Kokkos::TeamPolicy(numberOfElementsI, Kokkos::AUTO),
    KOKKOS_LAMBDA (const Kokkos::TeamPolicy<>::member_type& teamMember) {
        const int i = teamMember.team_rank();

        Kokkos::parallel_for(
            Kokkos::TeamThreadRange(teamMember, firstJ, lastJ),
            [=] (const int j) {
                Kokkos::parallel_for(
                    Kokkos::ThreadVectorRange(teamMember, firstK, lastK),
                    [=] (const int k) { /* ... */ }
                );
            }
        );
    }
);
\end{minted}

\paragraph{One-dimensional range}

\begin{minted}{c++}
Kokkos::TeamThreadRange range(teamMember, firstJ, lastJ);
Kokkos::ThreadVectorRange range(teamMember, firstK, lastK);
\end{minted}

\paragraph{Multi-dimensional range (dimension 2)}

\begin{minted}{c++}
Kokkos::TeamThreadMDRange<Kokkos::Rank<2>, Kokkos::TeamPolicy<>::member_type> range(teamMember, numberOfElementsJ, numberOfElementsK);
Kokkos::ThreadVectorMDRange<Kokkos::Rank<2>, Kokkos::TeamPolicy<>::member_type> range(teamMember, numberOfElementsL, numberOfElementsM);
\end{minted}

\section{Scratch memory}

Each team has access to a scratch memory pad, which has the team's lifetime, and is only accessible by the team's threads.

\subsection{Scratch memory space}

\begin{tblr}[theme=kokkostable]{lXl}
    Space level & Memory size & Access speed \\
    0 & Limited (tens of kilobytes) & Fast \\
    1 & Larger (few gigabytes) & Medium \\
\end{tblr}

Used when passing the team policy to the parallel construct and when creating the scratch memory pad.

\subsection{Create and populate}

\begin{minted}{c++}
// Define a scratch memory pad type
using ScratchPad = Kokkos::View<DataType, Kokkos::DefaultExecutionSpace::scratch_memory_space, Kokkos::MemoryTraits<Kokkos::Unmanaged>>;

// Compute how much scratch memory is needed (in bytes)
size_t bytes = ScratchPad::shmem_size(vectorSize);

// Create the team policy and specify the total scratch memory needed
Kokkos::parallel_for(
    "label",
    Kokkos::TeamPolicy<>(leagueSize, teamSize).set_scratch_size(spaceLevel, Kokkos::PerTeam(bytes)),
    KOKKOS_LAMBDA (const Kokkos::TeamPolicy<>::member_type& teamMember) {
        const int i = teamMember.league_rank();

        // Create the scratch pad
        ScratchPad scratch(teamMember.team_scratch(spaceLevel), vectorSize);

        // Initialize it
        Kokkos::parallel_for(
            Kokkos::TeamVectorRange(teamMember, vectorSize),
            [=] (const int j) { scratch(j) = getScratchData(i, j); }
        );

        // Synchronize
        teamMember.team_barrier();
    }
);
\end{minted}

\section{Atomics}

\subsection{Atomic operations}

\begin{tblr}[theme=kokkostable]{Xl}
    Operation & Replaces \\
    \mintinline{c++}{Kokkos::atomic_add(&x, y)} & \mintinline{c++}{x += y} \\
    \mintinline{c++}{Kokkos::atomic_and(&x, y)} & \mintinline{c++}{x &= y} \\
    \mintinline{c++}{Kokkos::atomic_dec(&x)} & \mintinline{c++}{x--} \\
    \mintinline{c++}{Kokkos::atomic_div(&x)} & \mintinline{c++}{x /= y} \\
    \mintinline{c++}{Kokkos::atomic_inc(&x)} & \mintinline{c++}{x++} \\
    \mintinline{c++}{Kokkos::atomic_lshift(&x, y)} & \mintinline{c++}{x = x << y} \\
    \mintinline{c++}{Kokkos::atomic_max(&x, y)} & \mintinline{c++}{x = std::max(x, y)} \\
    \mintinline{c++}{Kokkos::atomic_min(&x, y)} & \mintinline{c++}{x = std::min(x, y)} \\
    \mintinline{c++}{Kokkos::atomic_mod(&x, y)} & \mintinline{c++}{x }{\scriptsize\texttt{\%}}\mintinline{c++}{= y} \\
    \mintinline{c++}{Kokkos::atomic_mul(&x)} & \mintinline{c++}{x *= y} \\
    \mintinline{c++}{Kokkos::atomic_nand(&x, y)} & \mintinline{c++}{x = !(x && y)} \\
    \mintinline{c++}{Kokkos::atomic_or(&x, y)} & \mintinline{c++}{x |= y} \\
    \mintinline{c++}{Kokkos::atomic_rshift(&x, y)} & \mintinline{c++}{x = x >> y} \\
    \mintinline{c++}{Kokkos::atomic_sub(&x, y)} & \mintinline{c++}{x -= y} \\
    \mintinline{c++}{Kokkos::atomic_store(&x, y)} & \mintinline{c++}{x = y} \\
    \mintinline{c++}{Kokkos::atomic_xor(&x, y)} & \mintinline{c++}{x ^= y} \\
\end{tblr}

\subsection{Atomic exchanges}

\begin{tblr}[theme=kokkostable]{XX}
    Operation & Description \\
    \mintinline{c++}{Kokkos::atomic_exchange(&x, desired)} & Assign desired value to object and return old value \\
    \mintinline{c++}{Kokkos::atomic_compare_exchange(&x, expected, desired)} & Assign desired value to object if the object has the expected value and return the old value \\
\end{tblr}

\section{Mathematics}

\subsection{Math functions}

\begin{tblr}[theme=kokkostable]{Xl}
    Function type & List of functions (prefixed by \mintinline{c++}{Kokkos::}) \\
    Basic ops. & \mintinline{c++}{abs}, \mintinline{c++}{fabs}, \mintinline{c++}{fmod}, \mintinline{c++}{remainder}, \mintinline{c++}{fma}, \mintinline{c++}{fmax}, \mintinline{c++}{fmin}, \mintinline{c++}{fdim}, \mintinline{c++}{nan} \\
    Exponential & \mintinline{c++}{exp}, \mintinline{c++}{exp2}, \mintinline{c++}{expm1}, \mintinline{c++}{log}, \mintinline{c++}{log2}, \mintinline{c++}{log10}, \mintinline{c++}{log1p} \\
    Power & \mintinline{c++}{pow}, \mintinline{c++}{sqrt}, \mintinline{c++}{cbrt}, \mintinline{c++}{hypot} \\
    Trigonometric & \mintinline{c++}{sin}, \mintinline{c++}{cos}, \mintinline{c++}{tan}, \mintinline{c++}{asin}, \mintinline{c++}{acos}, \mintinline{c++}{atan}, \mintinline{c++}{atan2} \\
    Hyperbolic & \mintinline{c++}{sinh}, \mintinline{c++}{cosh}, \mintinline{c++}{tanh}, \mintinline{c++}{asinh}, \mintinline{c++}{acosh}, \mintinline{c++}{atanh} \\
    Error, gamma & \mintinline{c++}{erf}, \mintinline{c++}{erfc}, \mintinline{c++}{tgamma}, \mintinline{c++}{lgamma} \\
    Nearest & \mintinline{c++}{ceil}, \mintinline{c++}{floor}, \mintinline{c++}{trunc}, \mintinline{c++}{round}, \mintinline{c++}{nearbyint} \\
    Floating point & \mintinline{c++}{logb}, \mintinline{c++}{nextafter}, \mintinline{c++}{copysign} \\
    Comparisons & \mintinline{c++}{isfinite}, \mintinline{c++}{isinf}, \mintinline{c++}{isnan}, \mintinline{c++}{signbit} \\
\end{tblr}

Note that not all C++ standard math functions are available.

\subsection{Complex numbers}

\subsubsection{Create}

\begin{minted}{c++}
Kokkos::complex<double> complex(realPart, imagPart);
\end{minted}

\subsubsection{Manage}

\begin{tblr}[theme=kokkostable]{lX}
    Method & Description \\
    \mintinline{c++}{real()} & Returns or sets the real part \\
    \mintinline{c++}{imag()} & Returns or sets the imaginary part \\
\end{tblr}

\section{Utilities}

\subsection{Program interruption}

\begin{minted}{c++}
Kokkos::abort("message");
\end{minted}

\subsection{Print inside a kernel}

\begin{minted}{c++}
Kokkos::printf("format string", arg1, arg2);
\end{minted}

Similar to \mintinline{c++}{std::printf}.

\subsection{Timer}

\subsubsection{Create}

\begin{minted}{c++}
Kokkos::Timer timer;
\end{minted}

\subsubsection{Manage}

\begin{tblr}[theme=kokkostable]{lX}
    Method & Description \\
    \mintinline{c++}{seconds()} & Returns the time in seconds since construction or last reset \\
    \mintinline{c++}{reset()} & Resets the timer to zero \\
\end{tblr}

\subsection{Manage parallel environment}

\begin{tblr}[theme=kokkostable]{lX}
    Function & Description \\
    \mintinline{c++}{Kokkos::device_id()} & Returns the ID of the current device \\
    \mintinline{c++}{Kokkos::num_devices()} & Returns the number of devices available to the current execution space \\
\end{tblr}

\section{Macros}

\subsection{Essential macros}

\begin{tblr}[theme=kokkostable]{lX}
    Macro & Description \\
    \mintinline{c++}{KOKKOS_LAMBDA} & Replaces capture argument for lambdas \\
    \mintinline{c++}{KOKKOS_CLASS_LAMBDA} & Replaces capture argument for lambdas, captures \mintinline{c++}{this} \\
    \mintinline{c++}{KOKKOS_FUNCTION} & Functor attribute \\
    \mintinline{c++}{KOKKOS_INLINE_FUNCTION} & Inlined functor attribute \\
\end{tblr}

\subsection{Extra macros}

\begin{tblr}[theme=kokkostable]{lX}
    Macro & Description \\
    \mintinline{c++}{KOKKOS_VERSION} & Kokkos full version \\
    \mintinline{c++}{KOKKOS_VERSION_MAJOR} & Kokkos major version \\
    \mintinline{c++}{KOKKOS_VERSION_MINOR} & Kokkos minor version \\
    \mintinline{c++}{KOKKOS_VERSION_PATCH} & Kokkos patch level \\
    \mintinline{c++}{KOKKOS_ENABLE_*} & Any equivalent CMake option passed when building Kokkos, see installation cheat sheet \\
    \mintinline{c++}{KOKKOS_ARCH_*} & Any equivalent CMake option passed when building Kokkos, see installation cheat sheet \\
\end{tblr}

\end{multicols}

% add a watermark here

%\watermark

% add a notes section here

%\notessection

% add a notes section followed by a watermark here

%\notessection{\watermark}

\end{document}
